%\documentclass[]{article}
\documentclass[]{../../../../tex/classes/homework}
\usepackage{../../../../tex/packages/hebrewSupport}
\usepackage{../../../../tex/packages/mathShortcuts}
\usepackage{../../../../tex/packages/theoremsSupport}

%% Math packages
%\RequirePackage{ifxetex,ifluatex,amssymb,amsmath,mathrsfs,amsthm,witharrows,mathtools,mathdots,centernot}
%\RequirePackage{amsmath}
%\renewcommand{\qedsymbol}{$\blacksquare$} % end proofs with \blacksquare. Overwrites the defualts.
%\RequirePackage{bm}
%
%% Design
%\RequirePackage[margin=1in]{geometry}
%\RequirePackage[skip=4pt, indent=0pt]{parskip}
%\RequirePackage[normalem]{ulem}
%\renewcommand\labelitemi{$\bullet$}
%
%
%\newcommand\ff{\mathrm{f}}
%\newcommand\gcc{G^{\mathrm{SCC}}}
%\newcommand\squigT{\overset{T}{\squig}}
%\newcommand\squigd{\leftrightsquigarrow}
%\DeclareMathOperator{\low}{low}
%\DeclareMathOperator{\sol}{sol}
%\DeclareMathOperator{\val}{val}
%\DeclareMathOperator{\key}{key}
%\DeclareMathOperator{\Left}{left}
%\DeclareMathOperator{\Right}{right}
%
%
%\setlength{\parindent}{0pt}
%\setlength{\parskip}{4pt}

\DeclareMathOperator*{\argmin}{arg\,min}

\author{שחר פרץ}
\title{אלגוריתמים $\sim$ \textit{תרגיל בית 4}}
\begin{document}
	\maketitle
%	\ {} \\ \ \\ \ \\ 
	\section{}
	נתון גרף $G = (V, E)$ מכוון וצביעה של קשתות $c \co E \to \{\mathrm{brown}, \mathrm{purple}, \mathrm{green}, \mathrm{white}, \mathrm{red}\}$. נרצה לבדוק האם מתקיים:
	\[ 2 \cdot \#_{\mathrm{purple}}(C) > 3 \cdot \#_{\mathrm{brown}}(C) + 4  \cdot \#_{\mathrm{green}}(C) -5  \cdot \#_{\mathrm{white}}(C) \]
	כאשר $C$ הוא מעגל ב־$G$. 
	\begin{itemize}
		\item \textbf{אלגוריתם: }בהינתן $w \in E$, נגדיר את פונקציית המשקל הבאה על הגרף: 
		\[ w(e) = \begin{cases}
			-3 & c(e) = \mathrm{brown} \\
			-4 & c(e) = \mathrm{green} \\
			5 & c(w) = \mathrm{white} \\
			2 & c(w) = \mathrm{purple} \\
			0 & c(w) = \mathrm{red}
		\end{cases} \]
		נריץ את אלגוריתם BF על הגרף ונחזיר שכל המעגלים בעלי צביעה נאותה אם ורק אם BF החזיר שלא קיים מעגל שלילי. 
		\item \textbf{נכונות: }נבחין שמתקיים התנאי לעיל אמ''מ: 
		\[ 0 \le -3 \cdot \#_{\mathrm{brown}}(C) - 4  \cdot \#_{\mathrm{green}}(C) +5  \cdot \#_{\mathrm{white}}(C) + 2 \cdot \#_{\mathrm{purple}}(C) + 0 \cdot \#_{\mathrm{red}}(C) = w(C) \]
		כאשר השוויון ל־$w(C)$ מתקבל בקלות ממעבר על המעגל ושימוש בהגדרת $w$. סה''כ מעגל $C$ צבוע נאותה אמ''מ $w(C) \ge 0$. נסיק שכל המעגלים צבועים נאותה אמ''מ כל המעגלים מקיימים $w(C) \ge 0$, ומקונטראפוסיטיב זה שקול לכך שלא קיים מעגל $C$ שמקיים $w(C) < 0$, כלומר לא קיים מעגל בעל משקל שלילי. 
		\item \textbf{סיבוכיות: }כזמן ריצת הסיבוכיות של BF, שהיא $\bm{\oc(\sof E \sof V)}$. 
	\end{itemize}

%	\newpage \ \newpage 
	 \npage	
%	\ {} 
	 \section{}
	יהי גרף $G =(V, E)$ ופונקציית משקל $w \co E \to \R$. נניח שאין בגרף מעגלים שליליים. יהיו $s, t \in V$. נמצא את המסלול הקל ביותר שאיננו מק''ב (את קבוצת המסלולים שאינם מק''בים נסמן ב־$P$). 
%	\begin{itemize}
%		\item \textbf{אלגוריתם: }ראשית, נריץ BF רגיל כדי למצוא מה הוא ההילוך הקל ביותר, ונסמן את משקלו ב־$W$. לאחר מכן, נריץ BF עם שינוי קל לפעולת ה־\texttt{Relax}: אם מבצעים \en{\texttt{Relax($t$, $v$)}} \ (כאשר $v \in V$ שכן של $t$), לא נבצע עדכון במידה ו־$\dd[t] + w(t, v) = W$. 
%		\item \textbf{נכונות: }ראשית, נבחין ש־$P$ היא קבוצת כל ההילוכים $s \squig t$ שמשקלם אינו $W$. זאת כי מנכונות BF, $W$ הוא המשקל המינימלי. עכשיו נפנה להוכיח שהאלגוריתם מחזיר מסלול אמ''מ הוא מק''ב ב־$P$. 
%		\begin{itemize}
%			\item נניח שהאלגוריתם מחזיר מסלול. נראה שהוא מק''ב ב־$P$. הוא בבירור ב־$P$ כי האלגוריתם יחזיר את המסלול שבזמן ה־\texttt{Relax} משקלו ב־$t$ היה $\dd[t]$, ומהגדרת ה־\texttt{Relax} המשופר לא ייתכן $\dd[t] = W$. הוא גם כן מק''ב כי BF פועל כרגיל בעבור כל $v$ שכן של $t$ ללא שינויים (למעשה כאן משתמשים בהנחה שאין בגרף מעגלים שליליים) ולכן לאותם שכנים באיטרציה לפני ה־\texttt{Relax} האפקטיבי האחרון על $t$ שמור ב־$\dd[v]$ מסלולם הקצר ביותר ל־$s$, ואז מלמה 7 סיימנו. 
%			\item נניח שהאלגוריתם נכשל בלהחזיר מסלול. במקרה זה $\dd[t] = \infty$ (כלומר, לא ניתן לחזור אחורה ולשחזר מסלול). אם בגרף המקורי $G$ אין אף מסלול $s \squig t$, סיימנו. אחרת, נסיק שכל המסלולים המסתיימים ב־$t$ ממשקל $W$, ואז $P = \varnothing$, וגם סיימנו. 
%		\end{itemize}
%		\item \textbf{סיבוכיות: }ריצת BF הראשונה תארך $\oc(\sof V \sof E)$. השינוי ל־\texttt{Relax} בריצה השנייה לא משפיע על זמן הריצה כי מדובר בפעולה בעלת סיבוכיות קבועה, ומספר באיטרציות בלולאות ה־For של BF לא תשתנה. 
%	\end{itemize}
	\begin{itemize}
		\item \textbf{אלגוריתם: }נריץ BF על $s$ לפתרון SSSP על $s$ (למציאת $\dg(s, v)$ לכל $v \in V$), ובגרף המשוחלף $G^{T}$ נריץ BF על $t$ לפתרון SSSP (למציאת $\dg(v, t)$ לכל $v \in V$ בגרף המקורי). נבחין שבתהליך הזה מצאנו את $\dg(s, t)$ המק''ב. עתה נחפש את: 
		\[ \argmin_{(v, u) \in E} \underbrace{\dg(s, v) + w(v, u) + \dg(u, t)}_{: =\wg(v, u)} \neq \dg(s, t) \]
		כלומר, את ה־$(v, u) \in E$ הקשת שעבורה הערך $\dg(s, v) + w(v, u) + \dg(u, t)$ מינימלי אך לא שווה ל־$\dg(s, t)$, ונחזיר את המק''ב שהיא משרה (נוכל לשחזר באמצעות ה־$\pi$־ים מ־BF מסלולים $u \squig t$, מסלול $s \squig v$ ולהכניס קשת $(v, u)$ באמצע). במידה ו־$\dg(s, v) = \infty \lor \dg(u, t) = \infty$ לכל $(v, u)$ שעבורה $\wg(v, u) \neq \dg(s, t)$, נחזיר שלא קיים מסלול כזה ב־$P$. 
		\item \textbf{נכונות: }נוכיח שהאלגוריתם מחזיר מסלול אמ''מ הוא המינימלי מבין אלו שאינם ב־$P$. 
		\begin{itemize}
			\item אם האלגוריתם לא מחזיר אף מסלול, יהי מסלול $s \squig t$ ונראה שמשקלו $\dg(s, t)$. אם בשלילה משקלו אינו $\dg(s, t)$ אז בהכרח קיימת $(v, u)$ עליו כך ש־$\wg(v, u) \neq \dg(s, t)$, אחרת $\wg(v, u) = \dg(s,t)$ לכל $(v, u)$ ואז מאינדוקציה הוא בהכרח מק''ב. דהיינו משקלו $\dg(s, t)$ כנדרש. 
			\item נניח שהאלגוריתם מחזיר מסלול $p \co s \squig t$ ממשקל $x$, ויהי מסלול $s \squig t$ ב־$P$ אחר כלשהו. משום שהוא איננו מק''ב, מנימוקים שראינו קודם בהכרח קיימת עליו קשת $(v, u)$ כך ש־$\wg(v, u) \neq \dg(s, t)$ ולכן איננו קל ביותר. מהיות $x$ המינימלי בעבור פונ' המשקל $\wg$, ומשום שלכל מסלול $p'$ מתקיים בעבור קשת כלשהי $w(p') = \wg(v, u)$, בהכרח $p$ המסלול הקצר ביותר מבין כל המסלולים שאינם מק''בים, כנדרש. 
		\end{itemize}
		\item \textbf{סיבוכיות: }עלות הרצת BF פעמיים תהיה $\oc(\sof V \sof E)$. עלות חישוב הגרף המשוחלף $\oc(\sof E)$ וכנ''ל לגבי מציאת ה־$(v, u)$ המינימלי. עלות מציאת המסלול $s \squig t$ תיארך במקרה הגרוע $\oc(\sof E)$ גם כן (כי במקרה הגרוע המסלול יעבור בכל הגרף). סה''כ $\bm{\oc(\sof V \sof E)}$. 
	\end{itemize}
	
%	\newpage \ \newpage 
	\npage
%	\ {} 
	 \section{}
	נתון $G = (V, E)$ גרף אציקלי מכוון, ושני צמתים $s, t \in V$. נניח שיש צביעה $c \co E \to \{0, 1\}$ ($0$ יקרא \textit{אדום} ו־$1$ יקרא \textit{כחול}). נסמן ב־$\bg(p)$ את ה\textit{הטיה} של מסלול $p$ שהיא מוגדרת להיות ההפרש בין מספר הקשתות הכחולות לאדומות על $p$. 
	\begin{enumerate}[(A)]
		\item נמצא אלגוריתם שמחזיר מסלול מ–$s$ ל־$t$ עם הטיה מקסימלית. 
		\begin{itemize}
			\item \textbf{אלגוריתם: }נגדיר את פונקציית המשקל: 
			\[ w_1(e) = (-1)^{c(e)} = \begin{cases}
				1 & c(e) = 0 \equiv \mathrm{red} \\
				-1 & c(e) = 1 \equiv \mathrm{blue}
			\end{cases} \]
			באופן דומה נגדיר את $w_2(e)\co E \to \R$ להיות $w_2 = -w_1(e)$. 
			נריץ את האלגוריתם מההרצאה למציאת מק''בים בגרף ב־DAG (מיון טופולוגי ואז \texttt{Relax} לפי הסדר) הן בעבור $w_1$ והן בעבור $w_2$ ונחזיר את המקסימום מבין משקל המק''ב לפי $w_1$ ומשקל המק''ב לפי $w_2$. 
			\item \textbf{נכונות: }נגדיר $\bg_1(p)$ להיות ההפרש בין הקשתות הכחולות לאדומות, ו־$\bg_2(p)$ ההפרש בין הקשתות האדומות לכחולות. מהגדרתם, $\bg_1(p) = -\bg_2(p)$. 
			
			יהי $p$ מסלול, נוכיח ש־$w_1(p)$ ממינימלי אמ''מ $\bg_1(p)$ מקסימלי. נבחין כי: 
			\[ \bg_1(p) = \#_{\mathrm{blue}}(p) - \#_{\mathrm{red}}(p) = -\cl{(-1)\cdot \#_{\mathrm{blue}}(p) + 1 \cdot \#_{\mathrm{red}}(p)} = -w(p) \]
			וסה''כ מצאנו את ערך $\bg_1$ המקסימלי. 
			באופן דומה מציאת מק''ב לפי $w_2$ מוצאת את ערך $\bg_2$ המקסימלי. סה''כ, מהגדרת ערך מוחלט, ההטייה המקסימלית תהיה המקסימום מבין $\bg_1(p)$ ו־$\bg_2(p)$, שכן: 
			\[ \bg(p) = \sof{\bg_1(p)} = \max\{\bg_1(p), - \bg_1(p)\} = \max\{\bg_1(p), \bg_2(p)\} \]
			ומשום שהאלגוריתם מחזיר את הגבוה מבינהם, סיימנו. 
			
			
			\item \textbf{סיבוכיות: }עלות הגדרת פונקציית המשקל היא $\oc(\sof E)$. עלות הרצת האלגוריתם מההרצאה היא $\oc(\sof E + \sof V)$. סה''כ סיבוכיות לינארית של $\oc(\sof E + \sof V)$ כנדרש. 
		\end{itemize}
		\item נמצא אלגוריתם שמוצא מסלול $s \squig t$ עם הטיה מינימלית. 
		\begin{itemize}
			\item \textbf{אלגוריתם: }נבצע את אותו האלגוריתם כמו קודם לכן, אך במקום למצוא מסלול קל ביותר נמצא מסלול כבר ביותר. נעשה זאת ע''י מיון טופולוגי של איברי הגרף, ולכן מכן בעת מעבר על $v$ לכל $u \in \adj v$ (כלומר $u$ שכן של $v$) נגדיר: 
			\[ \dd[u] = \max\{\dd[v], \dd[v] + w(v, u)\} \]
			ונטען שבסוף הריצה $\dd[v] = \dg(s, v)$ (נאתחל $\dd[s] = 0$ ו־$\dd[v] = \infty$ לכל $v \neq s$) (\textit{הערה: }לשימושנו בשאלה זו $\dg(s, v)$ המסלול \textbf{הכבד} ביותר מ־$s$ ל־$v$). לבסוף נחזיר את המקסימום מבין הערך המוחלט של משקל המק''בים לפי $w_1$ ולפי $w_2$. 
			\item \textbf{נכונות: }יש להראות שהאלגוריתם מוצא מסלול כבד ביותר במקום קל ביותר. לאחר מכן, קל הטיעונים לעיל עדיין יהיו תקפים. 
			
			נסמן ב־$P_v$ את קבוצת הקודקודים שלהם קשת הנכנסת ל־$v$. בבירור: 
			\[ \dg(s, v) = \max_{u \in P_v}\dg(s, u) + w(v, u) \]
			שכן המסלול הכבד ביותר $s \squig p \to v$ (ייתכן $p = s$) הוא ממשקל $\dg(s, p) + v$ מהגדרה, ולכן רק נותר לקחת את המקסימום מבין כל ה־$p$־ים האפשריים. באינדוקציה ומהגדרת סדר מיון טופולוגי, כאשר נסיים לעבוד על הצומת ה־$u$ ונמשיך ממנו הלאה, כבר עדכנו את $\dd[u]$ בהתאם לכל הקודקודים שנכנסים אליו. לכן $\dg(s, v) = \dd[v]$ בסיום הריצה. 
			
			עדיין מתרחש $\bg(p) = \max\{\bg_1(p), \bg_2(p)\}$, ולכן לאחר המסלול $p_1$ המינימלי בעבור $\bg_1(p_1)$ והמינימלי בעבור $\bg_2(p_2)$, נותר להחזיר את המקסימום מבינהם ולקבל תשובה נכונה. 
			\item \textbf{סיבוכיות: }עלות מיון טופולוגי היא $\oc(\sof V + \sof E)$. האלגוריתם שלנו יעבור על כל $v \in V$ ועל כל שכיניו, ומלמת לחיצות הידיים סה''כ $\oc(\sof V + \sof E)$. נחזור על כל זה פעמיים, בעבור שתי פונקציות משקל שונות, ולכן סה''כ $\bm{\oc(\sof V + \sof E)}$ סיבוכיות לינארית. 
		\end{itemize}
	\end{enumerate}
	
%	\newpage \ \newpage 
	\npage
%	\ {} 
	\section{}
	נתון $G = (V, E)$, קודקוד $s \in V$ ו־$w \co E \to A$. נמצא פתרון ל־SSSP בזמן לינארי תחת התניות מסוימות. 
	\begin{itemize}
		\item נניח ש־$A = \range w = \{1, 2\}$. 
		\begin{itemize}
			\item \textbf{אלגוריתם: }נגדיר את הגרף $G' = (V', E')$ הבא: לכל $e \in E$, אם $w(e) = 1$ נשאיר את $e$ כמו שהיא, ואם $w(e) = 2$ נחליף את $e = (v, u)$ ב־$(v, v_{e})$ ו־$(v_e, u)$ כאשר $v_e$ הוא איזשהו קודקוד חדש של $G'$. עתה נריץ BFS על $G'$ ונטען שערכי ה־$\dd$ שה־BFS שמר על כל $v \in V$ הם אורכי המסלולים הקצרים ביותר מ־$s$. 
			\item \textbf{נכונות: }נבחין ש־$\din(v_e) = \dout(v_e) = 1$ לכל $e \in E$ המקיימת $w(e) = 2$. יהי $v \in V$ ויהי $p$ מסלול $s \squig v$. נוכיח ש־$w(p)$ הוא אורך המסלול הקצר ביותר ב־$G'$. נבחין שלכל $e = (x, y)$ המקיימת $w(e) = 2$ מתקיים שאורך המסלול הקצר ביותר $x \squig y$ ב־$G'$ הוא $2$ (כי $(v, v_e), (v_e, u)$ ב־$G'$ אך $e$ לא), וכן אם $w(e) = 1$ אורך המסלול הקצר ביותר ב־$G'$ הוא $1$ (כי במקרה זה $e$ ב־$G'$). לכן לאחר מעבר וסכימה על הקשתות ב־$p$ נקבל שאורך המסלול הקצר ביותר $s \squig v$ הוא המשקל $w(p)$, כי $w(p) = \#_1(p) + \#_2(p)$. 
			
			עתה מנכונות BFS סיימנו; BFS שומר ב־$\dd$ את המסלול הקצר ביותר ב־$G'$, לכל $v \in V$, והראנו שמסלול קצר ביותר זה שווה בערכו למסלול הקל ביותר ב־$G$ המקורי. 
		\end{itemize}
		\item נניח ש־$A = \range w = \{0, 1, 2\}$ ושהגרף איננו מכוון. 
		\begin{itemize}
			\item \textbf{אלגוריתם: }במידה ובין $u$ ל־$v$ וקיים מסלול של קשתות שכולן ממשקל אפס, נאמר ש־$v, u$ \textit{שקולים} ונסמן $v \sim u$ (זהו יחס שקילות כי הגרף לא מכוון). נחשב את הגרף $G/_\sim = (\tl V, \tl E)$ שיוגדר באופן הבא: $\tl V = V/_\sim$, ולכל קשת $(v, u) \in E$, אם $w(v, u) > 0$ נוסיף את הקשת $([v]_\sim, [e]_\sim)$ לגרף. עתה, נשתמש באלגוריתם הזהה לאלגוריתם מהסעיף הקודם (הוא אומנם על גרף מכוון, אבל אפשר להפוך כל קשת לקשת דו־כיוונית ובכך להשתמש בו) בין $[s]_\sim$ לכל קודקוד אחר. לכל $v \in [u]_\sim$, נגדיר $\dg(s, u) = \dg([s]_\sim, [u]_\sim)$ ונסיים. 
			\item \textbf{נכונות: }נוכיח שאכן לכל $v \in [u]_\sim$ מתקיים $\dg(s, u) = \dg([s]_\sim, [u]_\sim)$. ההוכחה פשוטה: כל קשת $([v_i]_\sim, [v_{i + 1}]_\sim) \in \tl E$ במסלול ב־$G/_\sim$ ניתנת להמרה לקשת $(u^{i}_i, u^{i}_{i +1})$ בעבור $u^{i}_i \in [v_i]_{\sim}, u^{i}_{i + 1} \in [v_{i + 1}]_\sim$, ומעבר בין $u^{i}_{i}$ ל־$u^{i + 1}_{i}$ תוכל להתבצע ע''י מעבר בקשתות ממשקל $0$ שלא ישפיעו על משקל המסלול. המרה זו חח''ע ולכן עובדת בשני הכיוונים. מנכונות סעיף א' מצאנו את המק''בים בגרף $G/_\sim$, ועתה הוכחנו שניתן להמירם למק''בים מאותו המשקל ב־$G$ (ולהפך), ובכך סיימנו את ההוכחה. 
		\end{itemize}
		\item נניח ש־$A = \{0 \dots k\}$. נסמן $2026 =: k$. 
		\begin{itemize}
			\item \textbf{אלגוריתם: }נשתמש ב־Dijkstra, אך נעשה שינוי לערמת המינימום: 
			\begin{itemize}[*]
				\item \textbf{תיאור מבנה הנתונים: }נתחזק מערך בגודל $k \sof V$ ובכל תא רשימה מקושרת של האיברים עם מפתח בעל ערך זה. נתחזק פוינטר למינימום (בהתחלה יצביע על 0). 
				\item \textbf{Insert: }נכניס לתא ה־$i$ את האיבר עם המפתח $i$. 
				\item \textbf{Delete-Min: }לפני הוצאת המינימום, נבצע עדכון למצביע: אם בתא (ממנו בתחילת הלולאה הוצאנו את המינימום) יש עוד דברים, נשאר במקום. אחרת, נעבור לינארית על התאים עד שנגיע לתא הלא ריק הבא. נשמיד משהו שנמצא בתא שמצביע עליו מצביע המינימום. 
				\item \textbf{Dec-Key: }כדי להוריד $k$ מאובייקט ממשקל $i$ פשוט נעביר אותו לתא ה־$k - i$ ונוסיף אותו לרשימה המקושרת בתא. לא נעדכן את הפוינטר למינימום. 
			\end{itemize}
			\item \textbf{נכונות: }ישירות מנכונות דייקסטרא, אחת שנוכיח שמבנה הנתונים עובד. ברור שהאיבר בעל המפתח ה־$i$ יהיה במיקום ה־$i$, והחלק הלא־טרוויאלי היחיד הוא מדוע הפוינטר אכן מצביע למינימום לפני מחיקת המינימום והחזרתו. 
			\begin{itemize}
				\item לאחר ה־Insert: ה־Insert מתבצע רק בהתחלה, ובהתחלה הפוינטר של המינימום ל־$0$ (שם מופיע קודקוד ההתחלה). 
				\item לאחר Dec-Key: הנכונות כאן נובעת ספציפית מאופן הפעולה של דייקסטרא. Dec-Key מבוצע לאחר Relax אפקטיבי בין קודקוד $v$ שהוא געת המינימום בערמה לבין קודקוד $u$ שנפחית את משקלו. משום שאין לנו קשתות עם משקל שלילי, לאחר ה־Relax האפקטיבי מפתח $u$ יהיה $\dg(s, u) = \dg(s, v) + w(s, v) \ge \dg(s, v)$, וידוע ש־$\dg(s,v)$ המינימום, לכן גם לאחר ה־Dec-Key לא הקטנו את $u$ מתחת למינימום. 
				\item לפני ה־Delete-Min (החלק החשוב): כאשר הסתכלנו על ה־Insert ו־Dec-Key,  הראנו שהמינימום מעולם לא קטן בין קריאות ל־Del-Min, ולכן מעבר לינארי במעלה הערמה ימצא לנו את המינימום הבא כנדרש. 
			\end{itemize}
			\item \textbf{סיבוכיות: }$\bm{\oc\cl{\sof V + \sof E}}$, כי זו עלות הפעולות בדייקסטרא, וכל הפעולות יארכו $\oc(1)$ פרט ל־Delete-Min לשאורך כל הריצה יעבור על הערמה בגודל $\oc(k\sof V) = \oc(\sof V)$ (כי $k = 2026$ קבוע). 
		\end{itemize}
	\end{itemize}
	
%	\newpage \ \newpage 
	\npage
%	\ {} 
	\section{}
	נתון גרף $G = (V, E)$ ופונקציית משקל $w \co E \to \R$. תהי $A \subset V$ קבוצה ו־$s \in A$. ידוע שאם משקל קשת $w(u, v) < 0$, אז בהכרח $u \in A \land v \in V \setminus A$, וכן לא קיימת קשת $(u, v)$ שעבור $v \in V \setminus A \land V \in A$. 
	
	\begin{itemize}
		\item \textbf{אלגוריתם: }מבין הקשתות $A \to V \setminus A$ נמצא את הקשת המינימלית, נסמן את משקלה בערך מוחלט ב־$m$. עתה, נגדיר את פונקציית הפוטנציאל הבאה: 
		\[ \phi(x) = \begin{cases}
			0 & x \in A \\
			m & x \in V \setminus  A
		\end{cases} \]
		שמשרה את פונקציית המשקל $\wg(v, u) = w(v, u) + \phi(v) - \phi(u)$. עתה נפעיל דייקסטרא על הגרף החדש ונחזיר את התוצאה. 
		\item \textbf{נכונות: }נבחין ש־$\wg(v, u) \ge 0$. תהי $(v, u) \in E$
		\begin{itemize}
			\item אם $w(v, u) < 0$. נתון $u \in A \land v \in V \setminus A$. אז מהגדרת $\phi$ בבירור $\phi(u) = m \land \phi(v) = m$. לכן ואז $\wg(v, u) = w(v, u) + 0 + m$ ומשום ש־$m$ המשקל השלילי על הגודל המקסימלי סה''כ $\wg(v, u) > 0$. 
			\item אם $w(v, u) \ge 0$ אז מהנתון בהכרח $v, u \in A$ או $v, u \in V \setminus A$, ואז $\phi(v) = \phi(u)$ כלומר $\wg(v, u) = w(v, u)$. 
		\end{itemize}
		הראנו בהרצאה שכל פונקציית פוטנציאל $\phi(x) \ge 0$ משמרת מק''בים משני הכיוונים כלומר הרצת דייקסטרא תחזיר את התשובה הנכונה, כנדרש. 
		\item \textbf{סיבוכיות: }של דייקסטרא, $\bm{\oc(\sof E + \sof V \log \sof V)}$ (מציאת $m$ תיארך זמן של $\oc(\sof E)$ ולכן זניחה). 
	\end{itemize}
	
%	\newpage \ \newpage 
	\npage
%	\ {} 
	\section{}
	נתון $G = (V, E)$ עם צביעה לאדום וכחול של הקשתות. נמצא אלגוריתם בסיבוכיות $\oc(\sof V \sof E)$ שימצא לכל $s, t \in V$ מסלול $s \squig t$ שמספר שינויי הצבע לאורכו מינימלי. 
	
	\begin{itemize}
		\item \textbf{אלגוריתם: }נגדיר גרף חדש. לכל $v \in V$, שני צמתים חדשים: $v_R, v_B$. נוסף על כך, לכל $(v, u) \in E$, אם $v, u$ מאותו הצבע $C \in \{R, B\}$ נוסיף קשת $(v_C, u_C)$ עם משקל $0$, ואם $v, u$ מצבעים שונים $C \in \{R, B\}$ ו־$\bar C \in \{R, B\}$ בהתאמה, נוסיף קשת $(v_C, u_{\bar C})$ במשקל $1$. נסמן את הגרף שקיבלנו ב־$\tl G = \cl{\tl V, \tl E}$. 
		
		נבחין שקיבלנו גרף מכוון עם פונקציית משקל $w \co \tl V \to \{0, 1\}$. לכן, לכל $v_R \in V$, נוכל להפעיל את סעיף ג' עם $k = 2$ במקום $k = 2026$ לקבלת מסלולים קצרים ביותר מ־$v_R$ לשאר צמתי הגרף. נוכל לבצע אותו התהליך בעבור $v_B$. סה''כ לכל $s, t \in V'$ מצאנו מסלול קצר ביותר ב־$\tl G$ ואת משקלו. נטען שמספר שינויי הצבע המינמלי בין $s \in V$ ל־$t \in V$ הוא: 
		\[ \min\ccb{\dg_{\tl G}(s_R, t_R),\ \dg_{\tl G}(s_R, t_B),\ \dg_{\tl G}(s_B, t_R),\ \dg_{\tl G}(s_B, t_B)} \]
		
		\textit{הערה: }נוכל לדבר על מק''ב $s \squig t$ ב־$\tl G$ משום שלמרות שב־$\tl G$ כל קודקוד $v$ משתכפל לכדי $v_R, v_B$, יש זוג קשתות $v_R \lra s_B$ ולכן נוכל לתרגם מסלול $s_R \squig t_R$ למסלול $s_B \squig t_R$ או $s_B \squig t_B$, ולהפך. 
		\item \textbf{נכונות: }מנכונות סעיף ג', אכן נמצא את המסלול הקצר ביותר ב־$\tl G$. נוכיח שערך המק''ב $s \squig t$ ב־$\tl G$ הוא מספר שינויי הצבע המינימלי במסלול $s \squig t$. 
		לשם כך, נניח שמספר שינויי הצבע המינימלי במסלול $s \squig t$ ב־$G$ הוא $k$. נוכיח שהמק''ב $s_R \squig t_R$ ב־$\tl G$ הוא ממשקל $k$. לשם כך, נציג הליך לתרגום מסלולים ב־$G$ למסלולים ב־$\tl G$: 
		
			 נניח שיש מסלול $s \squig t$ ב־$G$ עם $k$ שינויי צבע, נסמנם בקשתות $(v_i, u_i)$, ונקבל (בה''כ נניח ששינוי הצבע הראשון והאחרון מאדום לכחול. התהליך יהיה זהה אחרת): 
			\[ s \squig v_1 \to u_1 \squig \cdots \squig v_k \to u_k \squig t \]
			מסלול ב־$G$ (ייתכן $u_k = t$ או $s = v_1$). נוכל לתרגמו למסלול ב־$\tl G$: 
			\[ s_R \squig_R {v_1}_R \to {u_1}_{B} \squig_B {v_2}_B \to {v_3}_R \squig \cdots \squig {v_k}_R \to {u_k}_R \squig_R t_R \]
			(כאשר $\squig_R$ מסמל מסלול אדום ו־$\squig_B$ מסלול כחול). נבחין שמשקל מסלול זה הוא $k$, שכן הקשתות היחידות בעלות משקל שאיננו $0$ הן הקשתות שעוברות בין $v_i$ ל־$u_i$ (ויש $k$ כאלו). 
			
			עוד נבחין שההליך שתיארנו להמרה בין מסלול ב־$G$ עם מספר מינימלי של שינויי צבע למק''ב ב־$\tl G$ ממיר אותנו ממסלול כלשהו ב־$G$ למסלול למינימלי $\dg(s_R, t_R),\ \dg(s_R, t_B),\ \dg(s_B, t_R),\ \dg(s_B, t_B)$, ולהפך נוכל להמיר את המינימלי מביניהם למסלול ב־$G$ (ההליך לעיל הפיך). זה מסיים את הוכחת הנכונות. 
		\item \textbf{סיבוכיות: }עלות בניית $\tl G$ תהיה $2 \sof V + \sof E$. עלות הרצת דייסטקרא משופר (של סעיף 4ג) לכל קודקוד בנפרד פעמיים תהיה: 
		\[ \oc\cl{2\sof V \cl{\sof V + \sof E}} = \oc\cl{\sof {V}^{2} + \sof V \sof E} = \bm{\oc\cl{{\sof V \sof E}}} \]
		נוכל להניח ש־$\sof V \le \sof E$ כי אחרת נוכל לפרק לרכיבי קשירות ולהפעיל את האלגוריתם לכל רכיב קשירות בנפרד. 
	\end{itemize}
	
%	\newpage \ \newpage 
	\npage
%	\ {} 
	\section{}
	נבצע שינוי באלגוריתם למציאת מק''בים המבוסס על כפל מטריצות כדי להפוך אותו לאלגוריתם למציאת אורך המעגל השלילי הקצר ביותר בגרף. 
	\begin{itemize}
		\item \textbf{אלגוריתם: }נוסיף מסלולים ממשקל $0$ בין כל צומת לעצמו. ניזכר בכפל מטריצות min-plus שהגדרנו בהרצאה. נחשב את $D^{(\ml)}$ כפי שהוגדרה בהרצאה לכל $\ml \in \{1 \dots n + 1\}$. נבצע חיפוש בינארי: נחשב באמצעות iterative squaring את החזקות $D^{(\ml)}$ עבור $\ml = 2^{k} \le n$, ונבדוק האם יש איברים שליליים על האלכסון ב־$D^{\floor{\frac{n}{2}}}$. נמשיך כך בחיפוש בינארי עד שנמצא את ה־$\ml$ המינימלי שעל אלכסונו נמצא איזשהו $D^{(\ml)}_{ii} < 0$ שלילי. החיפוש הבינארי יתבצע על $\ml \in \{1 \dots n + 1\}$. 
		\item \textbf{נכונות: }החישוב של $D^{(\ml)}$ נכון ושומר על השמורה ($d_{ij}$ המק''ב בין $i$ ל־$j$) לפי מה שהוכחנו בהרצאה. משום שאנו שואפים למצוא מעגל שלילי, אכן $D^{(\ml)}_{ii}$ יהיה שלילי אמ''מ אכן יש מעגל שלילי שמתחיל ומסתיים ב־$i$ באורך לכל היותר $\ml$. אנו מסתמכים כאן על הקשתות במשקל $0$ שהוספנו, שמבטיחות לנו שאם יש מעגל באורך $k$ שמתחיל ומסתיים ב־$i$, יש מעגל באורך $\ml$ לכל $\ml > k$ (זה גם מה שמאפשר לנו לבצע את החיפוש הבינארי). 
		
		נבחין שבמידה ויש מעגל שלילי, ישנו מעגל שלילי פשוט באורך $\ml = n + 1$ לכל היותר (גודל מעגל פשוט מקסימלי בגרף מכוון, וכל מעגל לא פשוט אפשר להפוך למעגל פשוט), ולכן האלגוריתם בהכרח ימצא אותו. 
		\item \textbf{סיבוכיות: }עלות הכפלת מטריצות תהיה $\oc(n^{3})$. החיפוש על האלכסון בכל איטרציה יארך $\oc(n)$. נבחין שבאיטרציה ה־$k$ של החיפוש הבינארי, במקרה הגרוע, נצטרך לעשות iterative-squaring על טווח מגודל $\frac{n}{2^{k}}$, דבר שיארך $\log\cl{\frac{n}{2^{k}}}$. יהיו $\log n$ הרצות של החיפוש הבינארי, וסה''כ נכפול מטריצות: 
		\[ \sum_{k = 1}^{\log n}\log\cl{\frac{n}{2^{k}}} = \sum_{i = 1}^{\log n}\cl{\log n - \log k} = \frac{\log ^{2}n}{2} = \oc\cl{\log^{2}n} \]
		ולכן סה''כ סיבוכיות האלגוריתם תהיה: 
		\[ \oc\cl{\log^{2}n\cl{n^{3} + n}} = \bm{\oc\cl{n^{3}\log^{2}n}} \]
	\end{itemize}
	
%	\newpage \ \newpage 
	\npage
%	\ {} 
	\section{}
	יהי גרף $G =(V, E)$ מכוון ופונקציית משקל $w \co E \to \R$. כמו כן יהיו $U_1, U_2 \subset V$ זרות. נמצא את המסלול המסודר הקל ביותר לכל זוג צמתים. 
	\begin{itemize}
		\item \textbf{אלגוריתם: }נבנה את הגרף $\tl G = (\tl V, \tl E)$ באופן הבא:
		\begin{itemize}
			\item נסמן $U = V \setminus (U_1 \uplus U_2)$ (נבחין ש־$U, U_1, U_2$ מכסות את $V$). 
			\item כל קודקוד $v \in U$ יומר לשני קודקודים $v_1, v_2$ ב־$\tl V$ (נסמן $V_1$ את מי שמכיל את הקודקודים מהצורה $v_1$ וב־$V_2$ את הקבוצה שמכילה את הקודקודים מהצורה $v_2$). את יתר הקודקודים ב־$U_1 \uplus U_2$ נכניס ישירות ל־$\tl V$. 
			\item לכל קשת $(v, u) \in E$: 
			\begin{itemize}[*]
				\item אם $u \in U_1 \land v \in U_2$, נתעלם מהקשת הזו ולא נכניס אותה ל־$\tl E$ (לא ייתכן שהיא תופיע במסלול מסודר). 
				\item אם $v \in U_1 \land u \in U_2$, נכניס אותה ישירות ל־$\tl E$. 
				\item אם $u, v \in U_i$ (בעבור $i \in\{1, 2\}$ קבוע) אז נכניס ישירות את $(u, v)$ ל־$\tl E$. 
				\item אם $u \in U_i \land v \in U$ אז נוסיף את הקשת $(v_i, u) \in V_1 \times U_i$. 
				\item אם $v \in U_i \land u \in U$ אז נוסיף את הקשת $(v, u_i) \in U_i \times V_2$. 
			\end{itemize}
			משקל כל הקשתות לעיל יהיה כמשקל הקשתות המקוריות. 
			נסכם: $V = U \cup U_1 \cup U_2$ (למעשה איחוד זר) ו־$\tl V = V_1 \cup V_2 \cup U_1 \cup U_2$ (למעשה גם כאן איחוד זר). 
			\item לכל קודקוד $v \in U$, נוסיף קשת $(v_1, v_2)$ במשקל אפס. 
		\end{itemize}
		עתה נריץ את האלגוריתם של ג'ונסון על הגרף $\tl G$. לכל $s, t \in V$: 
		\begin{itemize}
			\item אם $s \in U$ נחזיר את המסלול המתחיל ב־$s_1$. אחרת $s \in U_1 \cup U_2$ ונחזיר פשוט את המסלול המתחיל ב־$s$. 
			\item אם $t \in U$ נחזיר את המסלול הנגמר ב־$t_2$. אחרת $t \in U_1 \cup U_2$ ונחזיר פשוט את המסלול הנגמר ב־$t$. 
		\end{itemize}
		אם המשקל אינסופי, שקילות לכך שהמסלול לא קיים. 
		\item \textbf{נכונות: }הפלט בבירור תקין, שכן $V_1, V_2$ נגישים מ־$U_1$ אך $U_1$ נגיש רק דרך $V_1$ ו־$V_2$ נגיש רק דרך $U_2, V_1$. לכן המעבר בין $V_1$ ל־$V_2$ הוא יחיד (לאו דווקא קיים, ייתכן שהמסלול מתחיל מ־$U_2$) ומכאן גם שהמעבר בין $U_1$ ל־$U_2$ הוא יחיד.  
		
		עתה נראה שמסלול מסודר קל ביותר ב־$G$ אמ''מ הוא מק''ב ב־$\tl G$. נבחין שיש המרה חח''ע (אך לא על) טבעית בין מסלולים ב־$G$ למסלולים ב־$\tl G$: יהיו $v, u \in V$
		\begin{itemize}
			\item אם $u \in U_2$, כל קודקוד מ־$U$ שנתקל בו לאחר המעבר ב־$U_2$ יומר לקודקוד מ־$V_2$ ע''י $U \ni u \mapsto u_2 \in V_2$ (לא נוכל לחזור ל־$U_1$ לאחר ביקור ב־$U_2$ ולכן אין צורך להתעסק במקרה זה). 
			\item אם $u \in U$, ונניח שיש לנו מסלול מסודר. אם המסלול מוכל ב־$U \cup V_1$ אין בעיה שכן נוכל להמיר את $U \mapsto V_1$. אם וכאשר נעבור ל־$V_2$ נפעל כפי שמתואר בעבור מסלול המתחיל מ־$V_2$ (המסלול לא יחזור חזרה ל־$U_1$ ולכן נוכל לשרשרם). 
			\item אם $u \in U_1$ נפעל כפי שפעלנו כאשר $u \in U$. 
		\end{itemize}
		ההמרה הזו חח''ע ממסלולים מסודרים ב־$G$ למסלולים ב־$\tl G$. אך נבחין שנוכל להפוך אותה בעבור מסלולים שמתחילים מ־$V_1 \cup U_1 \cup U_2$ ב־$\tl G$ ולהמירם למסלולים (מסודרים, כפי שהראנו) ב־$G$. 
		
		ההמרה הזו בבירור משמרת משקל כי הקשתות ה''נוספות`` היחידות שנדרש לעבור עליהן הן קשתות המעבר מ־$V_1$ ל־$V_2$ והן ממשקל $0$. 
		
		סה''כ מסלול קל ביותר מבין המסודרים ב־$G$ אמ''מ הוא מק''ב ב־$\tl G$ (כאשר קודקודים מ־$U$ יומרו לקודקודים מ־$V_1$). 
		
		מכאן, הנכונות אוטומאטית מנכונות האלגוריתם של ג'ונסון. 
		\item \textbf{סיבוכיות: }עלות בניית הגרף $\tl G$ תהיה $\oc\cl{\sof V  + \sof E}$. עלות הרצת האלגוריתם של ג'ונסון תהיה $\oc({\sof V^{2}\log \sof V + \sof V \sof E})$. סה''כ $\bm{\oc\cl{\sof V^{2}\log \sof V + \sof V \sof E}}$. 
	\end{itemize}
	
	\ndoc
\end{document}